\documentclass[12pt]{article}
\usepackage[utf8]{inputenc}
\usepackage[brazilian]{babel}
\usepackage{amsmath,amssymb,amsfonts}
\usepackage{geometry}
\usepackage{bm}

\geometry{verbose,tmargin=2.5cm,bmargin=2.5cm,lmargin=2.5cm,rmargin=2.5cm}

\title{Dedução Formal do Modelo MPC Incremental com Estado Aumentado $[x; u_{\text{prev}}]$}
\author{Johannes Fischer}
\date{Julho de 2026}

\begin{document}

\maketitle

\section{Representação do Sistema Físico}

Considere um sistema linear genérico discreto no tempo com $n_x$ estados, $n_u$ entradas, $n_y$ saídas e $n_d$ perturbações, representado pelas seguintes equações:
\begin{subequations}
\label{eq:sistema_fisico}
\begin{align}
x(t+1) &= A x(t) + B u(t) \\
y(t) &= D u(t) + E d(t)
\end{align}
\end{subequations}
onde:
\begin{itemize}
    \item $x(t) \in \mathbb{R}^{n_x}$ é o vetor de estados físicos.
    \item $u(t) \in \mathbb{R}^{n_u}$ é o vetor de entradas físicas de controle.
    \item $d(t) \in \mathbb{R}^{n_d}$ é o vetor de perturbações externas.
    \item $y(t) \in \mathbb{R}^{n_y}$ é o vetor de saídas algébricas do sistema.
    \item $A \in \mathbb{R}^{n_x \times n_x}$, $B \in \mathbb{R}^{n_x \times n_u}$, $D \in \mathbb{R}^{n_y \times n_u}$ e $E \in \mathbb{R}^{n_y \times n_d}$ são as matrizes do modelo.
\end{itemize}

---

\section{Formulação Incremental e Aumento de Estados}

Para introduzir ação integradora no controlador, definimos o incremento do sinal de controle $\Delta u(t) \in \mathbb{R}^{n_u}$ como:
\begin{equation}
\label{eq:delta_u}
\Delta u(t) = u(t) - u(t-1) \implies u(t) = u(t-1) + \Delta u(t)
\end{equation}

Definimos um novo vetor de estados aumentado $\tilde{x}(t) \in \mathbb{R}^{n_x + n_u}$ que armazena os estados físicos atuais e as entradas de controle aplicadas no passo anterior:
\begin{equation}
\tilde{x}(t) = \begin{bmatrix} x(t) \\ u(t-1) \end{bmatrix}
\end{equation}

\subsection{Equação de Atualização de Estados}
Substituindo \eqref{eq:delta_u} na dinâmica física de estados \eqref{eq:sistema_fisico}a:
\begin{align}
x(t+1) &= A x(t) + B (u(t-1) + \Delta u(t)) \nonumber \\
       &= A x(t) + B u(t-1) + B \Delta u(t)
\end{align}

A dinâmica de atualização do sinal de controle passado é:
\begin{equation}
u(t) = I_{n_u} u(t-1) + I_{n_u} \Delta u(t)
\end{equation}

Agrupando ambas as equações na forma matricial aumentada:
\begin{equation}
\begin{bmatrix} x(t+1) \\ u(t) \end{bmatrix} = \begin{bmatrix} A & B \\ \bm{0}_{n_u \times n_x} & I_{n_u} \end{bmatrix} \begin{bmatrix} x(t) \\ u(t-1) \end{bmatrix} + \begin{bmatrix} B \\ I_{n_u} \end{bmatrix} \Delta u(t)
\end{equation}
que pode ser reescrita compactamente como:
\begin{equation}
\label{eq:estado_aumentado_atualizado}
\tilde{x}(t+1) = \tilde{A} \tilde{x}(t) + \tilde{B} \Delta u(t)
\end{equation}
onde:
\begin{equation}
\tilde{A} = \begin{bmatrix} A & B \\ \bm{0} & I \end{bmatrix}_{(n_x+n_u) \times (n_x+n_u)}, \quad \tilde{B} = \begin{bmatrix} B \\ I \end{bmatrix}_{(n_x+n_u) \times n_u}
\end{equation}

\subsection{Equação de Saída Aumentada}
A saída do sistema original \eqref{eq:sistema_fisico}b depende diretamente de $u(t)$. Substituindo $u(t) = u(t-1) + \Delta u(t)$ na equação de saída:
\begin{align}
y(t) &= D (u(t-1) + \Delta u(t)) + E d(t) \nonumber \\
     &= D u(t-1) + D \Delta u(t) + E d(t)
\end{align}

Escrevendo $y(t)$ em termos do estado aumentado $\tilde{x}(t)$:
\begin{equation}
\label{eq:saida_aumentada}
y(t) = \tilde{C} \tilde{x}(t) + D \Delta u(t) + E d(t)
\end{equation}
onde a matriz de observação aumentada é:
\begin{equation}
\tilde{C} = \begin{bmatrix} \bm{0}_{n_y \times n_x} & D \end{bmatrix}_{n_y \times (n_x+n_u)}
\end{equation}

\section{Predições Futuras ao Longo do Horizonte $N$}

Dadas as condições no instante $t$, deseja-se prever a trajetória futura dos estados e das saídas sobre um horizonte de predição $N$.

\subsection{Predição de Estados}
Aplicando a equação \eqref{eq:estado_aumentado_atualizado} de forma recursiva:
\begin{align}
\tilde{x}(t+1|t) &= \tilde{A} \tilde{x}(t) + \tilde{B} \Delta u(t) \nonumber \\
\tilde{x}(t+2|t) &= \tilde{A}^2 \tilde{x}(t) + \tilde{A}\tilde{B} \Delta u(t) + \tilde{B} \Delta u(t+1) \nonumber \\
&\vdots \nonumber \\
\tilde{x}(t+N|t) &= \tilde{A}^N \tilde{x}(t) + \sum_{i=1}^N \tilde{A}^{N-i} \tilde{B} \Delta u(t+i-1)
\end{align}

Definindo o vetor empilhado de estados futuros $\bm{\tilde{X}}$ e o vetor de incrementos de controle $\bm{\Delta U}$:
\begin{equation}
\bm{\tilde{X}} = \begin{bmatrix} \tilde{x}(t+1|t) \\ \vdots \\ \tilde{x}(t+N|t) \end{bmatrix}_{N(n_x+n_u) \times 1}, \quad
\bm{\Delta U} = \begin{bmatrix} \Delta u(t) \\ \vdots \\ \Delta u(t+N-1) \end{bmatrix}_{N n_u \times 1}
\end{equation}
escrevemos a trajetória futura de estados na forma compacta:
\begin{equation}
\label{eq:predicao_estados_matricial}
\bm{\tilde{X}} = F_x \tilde{x}(t) + G_x \bm{\Delta U}
\end{equation}
onde:
\begin{equation}
F_x = \begin{bmatrix} \tilde{A} \\ \tilde{A}^2 \\ \vdots \\ \tilde{A}^N \end{bmatrix}, \quad
G_x = \begin{bmatrix}
\tilde{B} & \bm{0} & \dots & \bm{0} \\
\tilde{A}\tilde{B} & \tilde{B} & \dots & \bm{0} \\
\vdots & \vdots & \ddots & \vdots \\
\tilde{A}^{N-1}\tilde{B} & \tilde{A}^{N-2}\tilde{B} & \dots & \tilde{B}
\end{bmatrix}
\end{equation}

\subsection{Predição de Saídas}
Para obter a trajetória de saídas futuras, calculamos a saída prevista $y(t+k|t)$ para cada passo do horizonte $k = 1, 2, \dots, N$ a partir da equação de saída aumentada \eqref{eq:saida_aumentada}:
\begin{equation}
y(t+k|t) = \tilde{C} \tilde{x}(t+k|t) + D \Delta u(t+k) + E d(t+k)
\end{equation}

Para compreender a estrutura matricial resultante, detalhamos os primeiros passos de predição:
\begin{itemize}
    \item {\bf Para $k=1$:}
    \begin{align}
    y(t+1|t) &= \tilde{C} \tilde{x}(t+1|t) + D \Delta u(t+1) + E d(t+1) \nonumber \\
             &= \tilde{C} \left[ \tilde{A} \tilde{x}(t) + \tilde{B} \Delta u(t) \right] + D \Delta u(t+1) + E d(t+1) \nonumber \\
             &= \tilde{C}\tilde{A} \tilde{x}(t) + \tilde{C}\tilde{B} \Delta u(t) + D \Delta u(t+1) + E d(t+1)
    \end{align}
    
    \item {\bf Para $k=2$:}
    \begin{align}
    y(t+2|t) &= \tilde{C} \tilde{x}(t+2|t) + D \Delta u(t+2) + E d(t+2) \nonumber \\
             &= \tilde{C} \left[ \tilde{A}^2 \tilde{x}(t) + \tilde{A}\tilde{B} \Delta u(t) + \tilde{B} \Delta u(t+1) \right] + D \Delta u(t+2) + E d(t+2) \nonumber \\
             &= \tilde{C}\tilde{A}^2 \tilde{x}(t) + \tilde{C}\tilde{A}\tilde{B} \Delta u(t) + \tilde{C}\tilde{B} \Delta u(t+1) + D \Delta u(t+2) + E d(t+2)
    \end{align}
\end{itemize}

Generalizando este processo recursivo para um passo genérico $k = j$:
\begin{equation}
y(t+j|t) = \tilde{C}\tilde{A}^j \tilde{x}(t) + \sum_{i=1}^j \tilde{C}\tilde{A}^{j-i}\tilde{B}\Delta u(t+i-1) + D\Delta u(t+j) + E d(t+j)
\end{equation}

Empilhamos as saídas futuras no vetor $\bm{Y} \in \mathbb{R}^{N n_y}$ e as perturbações futuras medidas no vetor $\bm{D} \in \mathbb{R}^{N n_d}$, definidos como:
\begin{equation}
\bm{Y} = \begin{bmatrix} y(t+1|t) \\ y(t+2|t) \\ \vdots \\ y(t+N|t) \end{bmatrix}_{N n_y \times 1}, \quad
\bm{D} = \begin{bmatrix} d(t+1) \\ d(t+2) \\ \vdots \\ d(t+N) \end{bmatrix}_{N n_d \times 1}
\end{equation}

A trajetória de saída pode ser expressa compactamente pela relação linear:
\begin{equation}
\label{eq:predicao_saida_matricial}
\bm{Y} = F_y \tilde{x}(t) + G_y \bm{\Delta U} + E_{pert} \bm{D}
\end{equation}
onde:
\begin{equation}
F_y = \begin{bmatrix} \tilde{C}\tilde{A} \\ \tilde{C}\tilde{A}^2 \\ \vdots \\ \tilde{C}\tilde{A}^N \end{bmatrix}_{N n_y \times (n_x+n_u)}, \quad
E_{pert} = \begin{bmatrix} 
E & \bm{0} & \dots & \bm{0} \\
\bm{0} & E & \dots & \bm{0} \\
\vdots & \vdots & \ddots & \vdots \\
\bm{0} & \bm{0} & \dots & E 
\end{bmatrix}_{N n_y \times N n_d}
\end{equation}

A matriz de resposta forçada de saídas $G_y \in \mathbb{R}^{N n_y \times N n_u}$ mapeia os incrementos futuros $\bm{\Delta U} = [\Delta u(t)^T \quad \dots \quad \Delta u(t+N-1)^T]^T$. Como o vetor $\bm{\Delta U}$ termina no passo $t+N-1$, o termo de acoplamento direto final $D\Delta u(t+N)$ é assumido nulo ($\Delta u(t+N) = \bm{0}$) pois está além do horizonte de controle projetado. Assim, agrupando os coeficientes das equações anteriores, obtemos a estrutura de $G_y$:
\begin{equation}
\label{eq:matriz_Gy_correta}
G_y = \begin{bmatrix}
\tilde{C}\tilde{B} & D & \bm{0} & \bm{0} & \dots & \bm{0} \\
\tilde{C}\tilde{A}\tilde{B} & \tilde{C}\tilde{B} & D & \bm{0} & \dots & \bm{0} \\
\tilde{C}\tilde{A}^2\tilde{B} & \tilde{C}\tilde{A}\tilde{B} & \tilde{C}\tilde{B} & D & \dots & \bm{0} \\
\vdots & \vdots & \vdots & \ddots & \ddots & \vdots \\
\tilde{C}\tilde{A}^{N-2}\tilde{B} & \tilde{C}\tilde{A}^{N-3}\tilde{B} & \tilde{C}\tilde{A}^{N-4}\tilde{B} & \dots & \tilde{C}\tilde{B} & D \\
\tilde{C}\tilde{A}^{N-1}\tilde{B} & \tilde{C}\tilde{A}^{N-2}\tilde{B} & \tilde{C}\tilde{A}^{N-3}\tilde{B} & \dots & \tilde{C}\tilde{A}\tilde{B} & \tilde{C}\tilde{B}
\end{bmatrix}
\end{equation}
Note que a matriz de acoplamento direto $D$ aparece deslocada na primeira super-diagonal da matriz $G_y$, pois o efeito instantâneo do incremento $\Delta u(t+k)$ afeta a saída $y(t+k|t)$ no mesmo passo, enquanto o efeito dinâmico do estado $\tilde{C}\tilde{B}$ atua com um atraso de um passo (multiplicando $\Delta u(t+k-1)$).

\section{Formulação Quadrática (QP) para o \texttt{quadprog}}

A função objetivo a ser minimizada busca conciliar o rastreamento das trajetórias de referência dos estados físicos e das saídas com a penalização do esforço de variação do controle, sendo expressa por:
\begin{equation}
\label{eq:custo_original}
J = (\bm{X}_f - \bm{X}_{\text{ref}})^T Q (\bm{X}_f - \bm{X}_{\text{ref}}) + (\bm{Y} - \bm{Y}_{\text{ref}})^T S (\bm{Y} - \bm{Y}_{\text{ref}}) + \bm{\Delta U}^T R \bm{\Delta U}
\end{equation}
onde:
\begin{itemize}
    \item $\bm{X}_f = [x(t+1|t)^T \quad \dots \quad x(t+N|t)^T]^T \in \mathbb{R}^{N n_x}$ representa a projeção futura apenas dos estados físicos reais.
    \item $\bm{X}_{\text{ref}} \in \mathbb{R}^{N n_x}$ é a trajetória de referência dos estados físicos (como os valores fixados de $SoC$ e $LOH$).
    \item $\bm{Y} \in \mathbb{R}^{N n_y}$ é a predição de saídas calculada em \eqref{eq:predicao_saida_matricial}.
    \item $\bm{Y}_{\text{ref}} \in \mathbb{R}^{N n_y}$ é a trajetória de referência das saídas (a demanda elétrica a ser atingida).
    \item $Q \in \mathbb{R}^{N n_x \times N n_x}, S \in \mathbb{R}^{N n_y \times N n_y}$ e $R \in \mathbb{R}^{N n_u \times N n_u}$ são as matrizes de ponderação do controle.
\end{itemize}

Nas seções a seguir, detalha-se a transição desta representação teórica para a forma quadrática padrão exigida para a programação computacional.

\subsection{Extração dos Estados Físicos da Matriz Aumentada}
\label{subsec:extracao}
O vetor de estados futuros predito pelo modelo aumentado em \eqref{eq:predicao_estados_matricial} é $\bm{\tilde{X}}$, que tem dimensão $N(n_x+n_u)$ pois inclui os controles passados $u(t+k-1)$ no vetor de estados. Como na função custo queremos penalizar apenas os desvios dos estados físicos $x(t+k)$, definimos uma matriz de seleção $T_{sel} \in \mathbb{R}^{N n_x \times N(n_x+n_u)}$ tal que:
\begin{equation}
\bm{X}_f = T_{sel} \bm{\tilde{X}} = T_{sel} F_x \tilde{x}(t) + T_{sel} G_x \bm{\Delta U}
\end{equation}

A matriz $T_{sel}$ é estruturada em blocos diagonais onde cada bloco seleciona os primeiros $n_x$ elementos de cada passo do horizonte:
\begin{equation}
T_{sel} = \begin{bmatrix} 
\begin{bmatrix} I_{n_x} & \bm{0}_{n_x \times n_u} \end{bmatrix} & \bm{0} & \dots & \bm{0} \\
\bm{0} & \begin{bmatrix} I_{n_x} & \bm{0}_{n_x \times n_u} \end{bmatrix} & \dots & \bm{0} \\
\vdots & \vdots & \ddots & \vdots \\
\bm{0} & \bm{0} & \dots & \begin{bmatrix} I_{n_x} & \bm{0}_{n_x \times n_u} \end{bmatrix}
\end{bmatrix}
\end{equation}

Definindo as matrizes de predição física $F_{xf} = T_{sel} F_x$ e $G_{xf} = T_{sel} G_x$, escrevemos:
\begin{equation}
\bm{X}_f = F_{xf} \tilde{x}(t) + G_{xf} \bm{\Delta U}
\end{equation}

\subsection{Expansão da Função Objetivo para Programação Quadrática}
\label{subsec:expansao}
Definimos os erros nominais (resposta livre sem novos incrementos menos as trajetórias de referência alvos) para estados e saídas como:
\begin{align}
\bm{e}_x &= F_{xf} \tilde{x}(t) - \bm{X}_{\text{ref}} \\
\bm{e}_y &= F_y \tilde{x}(t) + E_{pert} \bm{D} - \bm{Y}_{\text{ref}}
\end{align}

Com isso, reescrevemos as diferenças em relação às referências:
\begin{align}
\bm{X}_f - \bm{X}_{\text{ref}} &= \bm{e}_x + G_{xf} \bm{\Delta U} \\
\bm{Y} - \bm{Y}_{\text{ref}} &= \bm{e}_y + G_y \bm{\Delta U}
\end{align}

Substituindo estas expressões na função custo \eqref{eq:custo_original}:
\begin{align}
J &= (\bm{e}_x + G_{xf} \bm{\Delta U})^T Q (\bm{e}_x + G_{xf} \bm{\Delta U}) + (\bm{e}_y + G_y \bm{\Delta U})^T S (\bm{e}_y + G_y \bm{\Delta U}) + \bm{\Delta U}^T R \bm{\Delta U} \nonumber \\
  &= \bm{\Delta U}^T (G_{xf}^T Q G_{xf} + G_y^T S G_y + R) \bm{\Delta U} + 2 (\bm{e}_x^T Q G_{xf} + \bm{e}_y^T S G_y) \bm{\Delta U} + \text{constantes}
\end{align}

O solver `quadprog` do MATLAB minimiza a seguinte estrutura quadrática padrão:
\begin{equation}
\min_{\bm{\Delta U}} \frac{1}{2} \bm{\Delta U}^T H \bm{\Delta U} + f^T \bm{\Delta U}
\end{equation}

Por comparação direta de termos com a nossa função custo expandida, as matrizes $H$ e $f$ enviadas ao solver são definidas como:
\begin{align}
\label{eq:matrizes_qp_quadprog}
H &= 2 \left( G_{xf}^T Q G_{xf} + G_y^T S G_y + R \right) \\
f &= 2 \left( \bm{e}_x^T Q G_{xf} + \bm{e}_y^T S G_y \right)^T
\end{align}
*(Note que a constante de offset que não depende de $\bm{\Delta U}$ é omitida na otimização).*

\subsection{Propriedades de Otimização e Solução via QP}
\label{subsec:conclusao}
A formulação matemática proposta resulta em um problema de otimização quadrático convexo em relação à variável de decisão $\bm{\Delta U}$. Isso ocorre devido à linearidade das equações de predição dos estados e das saídas em relação aos incrementos de controle, associada à definição positiva da matriz de ponderação $R$. Portanto, as matrizes $H$ e $f$ obtidas em \eqref{eq:matrizes_qp_quadprog} definem de forma unívoca e rigorosa o problema a ser resolvido pelo solver \texttt{quadprog} do MATLAB.

\end{document}
